\section{Detección de Outliers}

Para detectar outliers lo que haremos será ocupar IQR. Por lo anterior notamos que la incidencia delictiva está fuertemente
sesgada a la derecha por la CDMX y Edomex, ya que disparan los valores frente a estados pequeños. El z-score asume normalidad
y los propios extremos inflan media y desviación, así que ocupar ésta medida reportaría resultados sesgados.


Vemos que del total de registros, 63128 fueron identificados como outliers, lo que representa aproximadamente el 15.25% del dataset. Esto es
un porcentaje considerable, pero dado el contexto de la delincuencia en México, es comprensible que existan registros con incidencias
extremadamente altas, especialmente en entidades como la CDMX y el EdoMex.

Sin embargo, es importante destacar que estos outliers no necesariamente representan errores en los datos, sino que reflejan la realidad de la
incidencia delictiva en ciertas regiones y periodos. Por lo tanto, en lugar de eliminar estos registros, se ha optado por marcarlos con una
nueva columna para su posterior análisis y consideración en cualquier modelado o interpretación de los datos. De esta manera, los datos se muestran
como sigue: